\section{Cosets and the Theorem of Lagrange}
\textbf{Conventions.} Permutations compose from right to left. The identity is $e$. A group generated by an element $a$ is denoted $\gen{a}$.

\exercise{10.1}{Cosets of $4\Z$ in $\Z$}
In additive notation, a coset is $a+4\Z=\{a+4k:k\in\Z\}$. Two representatives give the same coset precisely when their difference is divisible by $4$. The distinct cosets in $\Z$ are\[0+4\Z,\quad 1+4\Z,\quad 2+4\Z,\quad 3+4\Z.\]The division algorithm shows every allowed representative belongs to one of these, and their different residues show the listed cosets are distinct.

\exercise{10.2}{Cosets of $4\Z$ in $2\Z$}
In additive notation, a coset is $a+4\Z=\{a+4k:k\in\Z\}$. Two representatives give the same coset precisely when their difference is divisible by $4$. The distinct cosets in $2\Z$ are\[0+4\Z,\quad 2+4\Z.\]The division algorithm shows every allowed representative belongs to one of these, and their different residues show the listed cosets are distinct.

\exercise{10.3}{Cosets of $\gen{2}$ in $\Z_{12}$}
First list the subgroup generated by $2$ by repeatedly adding it modulo $12$:\[H=\gen{2}=\{0,2,4,6,8,10\}.\]Adding a representative to every element of $H$ gives all distinct cosets:\[\begin{aligned}0+H&=\{0,2,4,6,8,10\},\\
1+H&=\{1,3,5,7,9,11\}.\end{aligned}\]These sets are disjoint and together contain all $12$ residues, verifying that the list is complete.

\exercise{10.4}{Cosets of $\gen{4}$ in $\Z_{12}$}
First list the subgroup generated by $4$ by repeatedly adding it modulo $12$:\[H=\gen{4}=\{0,4,8\}.\]Adding a representative to every element of $H$ gives all distinct cosets:\[\begin{aligned}0+H&=\{0,4,8\},\\
1+H&=\{1,5,9\},\\
2+H&=\{2,6,10\},\\
3+H&=\{3,7,11\}.\end{aligned}\]These sets are disjoint and together contain all $12$ residues, verifying that the list is complete.

\exercise{10.5}{Cosets of $\gen{18}$ in $\Z_{36}$}
First list the subgroup generated by $18$ by repeatedly adding it modulo $36$:\[H=\gen{18}=\{0,18\}.\]Adding a representative to every element of $H$ gives all distinct cosets:\[\begin{aligned}0+H&=\{0,18\},\\
1+H&=\{1,19\},\\
2+H&=\{2,20\},\\
3+H&=\{3,21\},\\
4+H&=\{4,22\},\\
5+H&=\{5,23\},\\
6+H&=\{6,24\},\\
7+H&=\{7,25\},\\
8+H&=\{8,26\},\\
9+H&=\{9,27\},\\
10+H&=\{10,28\},\\
11+H&=\{11,29\},\\
12+H&=\{12,30\},\\
13+H&=\{13,31\},\\
14+H&=\{14,32\},\\
15+H&=\{15,33\},\\
16+H&=\{16,34\},\\
17+H&=\{17,35\}.\end{aligned}\]These sets are disjoint and together contain all $36$ residues, verifying that the list is complete.

\exercise{10.6}{Left cosets of a reflection subgroup in $D_4$}
Let $H=\{\rho_0,\mu_2\}$. For $a\in D_4$, the left coset is $aH=\{a,a\mu_2\}$. Reading the indicated products from Table 8.12 gives
\[\begin{aligned}\rho_0H&=\{\rho_0,\mu_2\},\\
\rho_1H&=\{\rho_1,\delta_2\},\\
\rho_2H&=\{\rho_2,\mu_1\},\\
\rho_3H&=\{\rho_3,\delta_1\}.\end{aligned}\]
The four pairs are disjoint and include all eight group elements, so these are all the left cosets.

\exercise{10.7}{Right cosets of the reflection subgroup}
A right coset is $H\rho_i=\{\rho_i,\mu_2\rho_i\}$, so the list is
\[\begin{aligned}H\rho_0&=\{\rho_0,\mu_2\},\\
H\rho_1&=\{\rho_1,\delta_1\},\\
H\rho_2&=\{\rho_2,\mu_1\},\\
H\rho_3&=\{\rho_3,\delta_2\}.\end{aligned}\]
These do not give the same partition as the left cosets: for example $\rho_1H=\{\rho_1,\delta_2\}$ while $H\rho_1=\{\rho_1,\delta_1\}$. Thus $H$ is not normal.

\exercise{10.8}{Reordered table for the reflection subgroup}
Order the rows and columns in consecutive left cosets from Exercise 6. The resulting table is
\[\begin{array}{c|cccccccc} & \rho_0 & \mu_2 & \rho_1 & \delta_2 & \rho_2 & \mu_1 & \rho_3 & \delta_1\\\hline
\rho_0 & \rho_0 & \mu_2 & \rho_1 & \delta_2 & \rho_2 & \mu_1 & \rho_3 & \delta_1\\
\mu_2 & \mu_2 & \rho_0 & \delta_1 & \rho_3 & \mu_1 & \rho_2 & \delta_2 & \rho_1\\
\rho_1 & \rho_1 & \delta_2 & \rho_2 & \mu_1 & \rho_3 & \delta_1 & \rho_0 & \mu_2\\
\delta_2 & \delta_2 & \rho_1 & \mu_2 & \rho_0 & \delta_1 & \rho_3 & \mu_1 & \rho_2\\
\rho_2 & \rho_2 & \mu_1 & \rho_3 & \delta_1 & \rho_0 & \mu_2 & \rho_1 & \delta_2\\
\mu_1 & \mu_1 & \rho_2 & \delta_2 & \rho_1 & \mu_2 & \rho_0 & \delta_1 & \rho_3\\
\rho_3 & \rho_3 & \delta_1 & \rho_0 & \mu_2 & \rho_1 & \delta_2 & \rho_2 & \mu_1\\
\delta_1 & \delta_1 & \rho_3 & \mu_1 & \rho_2 & \delta_2 & \rho_1 & \mu_2 & \rho_0\\\end{array}\]
The block $(\rho_0H)(\rho_1H)$ contains elements from different left cosets: $\rho_0\rho_1=\rho_1\in\rho_1H$, but $\mu_2\rho_1=\delta_1\in\rho_3H$. Both $\rho_0$ and $\mu_2$ represent the same coset $H$, so this demonstrates that changing a representative changes the product coset. Thus coset multiplication is not well defined. There is no coset group here, so it is isomorphic to neither $\Z_4$ nor $V$.

\exercise{10.9}{Left cosets of the half-turn subgroup}
Set $K=\{\rho_0,\rho_2\}$. The distinct left cosets are
\[\begin{gathered}
K=\{\rho_0,\rho_2\},\quad \rho_1K=\{\rho_1,\rho_3\},\\
\mu_1K=\{\mu_1,\mu_2\},\quad \delta_1K=\{\delta_1,\delta_2\}.
\end{gathered}\]
Each pair comes from multiplying the representative by $\rho_0$ and by $\rho_2$, and the four pairs partition the eight elements.

\exercise{10.10}{Right cosets of the half-turn subgroup}
The half-turn $\rho_2$ commutes with every square symmetry: algebraically, $r^2r^i=r^ir^2$ and $r^2(r^is)=(r^is)r^2$, using $sr^2=r^{-2}s=r^2s$. Hence $aK=Ka$ for each $a\in D_4$. The right cosets are therefore the same four pairs listed in Exercise 9; $K$ is normal.

\exercise{10.11}{Reordered table and the Klein quotient}
Using the order $\rho_0,\rho_2,\rho_1,\rho_3,\mu_1,\mu_2,\delta_1,\delta_2$, the table is
\[\begin{array}{c|cccccccc} & \rho_0 & \rho_2 & \rho_1 & \rho_3 & \mu_1 & \mu_2 & \delta_1 & \delta_2\\\hline
\rho_0 & \rho_0 & \rho_2 & \rho_1 & \rho_3 & \mu_1 & \mu_2 & \delta_1 & \delta_2\\
\rho_2 & \rho_2 & \rho_0 & \rho_3 & \rho_1 & \mu_2 & \mu_1 & \delta_2 & \delta_1\\
\rho_1 & \rho_1 & \rho_3 & \rho_2 & \rho_0 & \delta_1 & \delta_2 & \mu_2 & \mu_1\\
\rho_3 & \rho_3 & \rho_1 & \rho_0 & \rho_2 & \delta_2 & \delta_1 & \mu_1 & \mu_2\\
\mu_1 & \mu_1 & \mu_2 & \delta_2 & \delta_1 & \rho_0 & \rho_2 & \rho_3 & \rho_1\\
\mu_2 & \mu_2 & \mu_1 & \delta_1 & \delta_2 & \rho_2 & \rho_0 & \rho_1 & \rho_3\\
\delta_1 & \delta_1 & \delta_2 & \mu_1 & \mu_2 & \rho_1 & \rho_3 & \rho_0 & \rho_2\\
\delta_2 & \delta_2 & \delta_1 & \mu_2 & \mu_1 & \rho_3 & \rho_1 & \rho_2 & \rho_0\\\end{array}\]
Every $2\times2$ block lies entirely in a single coset, so multiplication descends to the four cosets. Label them $K,A=\rho_1K,B=\mu_1K,C=\delta_1K$. The induced table is
\[\begin{array}{c|cccc}&K&A&B&C\\\hline K&K&A&B&C\\A&A&K&C&B\\B&B&C&K&A\\C&C&B&A&K\end{array}\]
Every nonidentity element has order $2$. Thus the quotient is the Klein four-group $V$, not $\Z_4$.

\exercise{10.12}{Index of $\gen3$ in $\Z_{24}$}
The multiples of $3$ modulo $24$ are $0,3,6,9,12,15,18,21$, so the subgroup has order $8$. Lagrange's counting formula gives $(\Z_{24}:\gen3)=24/8=3$. Equivalently the cosets are the three residue classes modulo $3$.

\exercise{10.13}{Index of a reflection subgroup in $S_3$}
A reflection $\mu_1$ has order $2$, so $\gen{\mu_1}=\{e,\mu_1\}$. Since $|S_3|=6$, its index is $6/2=3$.

\exercise{10.14}{Index of a reflection subgroup in $D_4$}
The element $\mu_2$ is a reflection and has order $2$. Its subgroup has two elements, whereas $D_4$ has eight. The number of cosets is therefore $8/2=4$.

\exercise{10.15}{Index of the subgroup generated by $\sigma$}
The factors in $\sigma=(1,2,5,4)(2,3)$ overlap, so their orders cannot be combined by a least-common-multiple rule. Follow their images from right to left:
\[1\mapsto2\mapsto3\mapsto5\mapsto4\mapsto1.\]
Thus $\sigma=(1,2,3,5,4)$ has order $5$. Therefore $|\gen\sigma|=5$ and $(S_5:\gen\sigma)=5!/5=24$.

\exercise{10.16}{Index of the subgroup generated by $\mu$}
The cycles $(1,2,4,5)$ and $(3,6)$ are disjoint and have lengths $4,2$. Hence $\mu$ has order $\lcm(4,2)=4$. Its cyclic subgroup has four elements, and the index in $S_6$ is $6!/4=720/4=180$.

\exercise{10.17}{Definition of a left coset}
Require that $H$ be a \emph{subgroup} of $G$, not merely a subset, and that $a\in G$. Then the left coset of $H$ containing $a$ is $aH=\{ah:h\in H\}$. Because $e\in H$, it indeed contains $a=ae$. For an arbitrary subset, the proposed phrase ``containing $a$'' need not be true.

\exercise{10.18}{Definition of index}
The definition using the number of right cosets is correct. The number of left cosets is equal to it, so either convention defines the same index $(G:H)$. For finite $G$ this number is $|G|/|H|$; the definition itself does not require finiteness.

\exercise{10.19}{True or false}
\begin{enumerate}[label=(\alph*)]
\item True: $aH$ is defined for every subgroup $H\le G$ and every $a\in G$.
\item True: $|G|=(G:H)|H|$, so the number of cosets divides $|G|$.
\item True: a group of prime order is cyclic, and cyclic groups are abelian.
\item False: the finite subgroup $\{0\}$ of $\Z$ has singleton cosets $n+\{0\}$.
\item True: $H=eH$.
\item False: for example $\Z\le\R$ has cosets $x+\Z$.
\item True: for $n>1$, half the permutations are even, so $[S_n:A_n]=2$.
\item This is a statement of taste rather than a mathematical proposition; one may agree that the theorem is nice, but there is no objectively defined truth value to prove.
\item False: $V$ has order $4$ but no element of order $4$.
\item True: if $G=\gen g$ has order $m$ and $d\mid m$, then $g^{m/d}$ has order $d$.
\end{enumerate}

\exercise{10.20}{Different left and right partitions in an abelian group}
Impossible. For every $a\in G$ and $h\in H$, commutativity gives $ah=ha$, so $aH=Ha$. Thus the left and right coset partitions must agree for every subgroup of an abelian group.

\exercise{10.21}{A single coset}
Take $H=G$, for example $H=G=\Z_6$. Every $aH=aG=G$, so there is exactly one cell in the partition. Conversely an index-$1$ subgroup must be all of $G$, since its identity coset $H$ is the only cell.

\exercise{10.22}{Six cosets in a group of order six}
Take $G=\Z_6$ and the trivial subgroup $H=\{0\}$. Its six cosets are the six singletons, so the partition has six cells. The formula $[G:H]=6/1$ confirms the count.

\exercise{10.23}{Twelve cosets in a group of order six}
Impossible. Every coset is nonempty, so a partition of a six-element set has at most six cells. Equivalently, the index would force $|H|=6/12=1/2$, which is not a possible subgroup order.

\exercise{10.24}{Four cosets in a group of order six}
Impossible. Lagrange's formula would give $6=4|H|$, so $|H|=3/2$, not an integer. Thus the index of a subgroup of an order-$6$ group cannot be $4$.

\exercise{10.25}{Synopsis of Lagrange's theorem}
The left cosets of $H$ partition finite $G$, and each has $|H|$ elements by the bijection $h\mapsto ah$, so $|G|=[G:H]|H|$ and $|H|$ divides $|G|$.

\exercise{10.26}{The right-coset equivalence relation}
The relation is $a\sim_R b$ if and only if $ab^{-1}\in H$.
It is reflexive because $aa^{-1}=e\in H$.
If $ab^{-1}\in H$, inverse closure gives $(ab^{-1})^{-1}=ba^{-1}\in H$, so it is symmetric.
If $ab^{-1}\in H$ and $bc^{-1}\in H$, their product is $ac^{-1}\in H$, proving transitivity. Hence it is an equivalence relation; its class containing $a$ is $Ha$.

\exercise{10.27}{A bijection with a right coset}
Define $f:H\to Hg$ by $f(h)=hg$. Its values lie in $Hg$ by definition. If $f(h_1)=f(h_2)$, then $h_1g=h_2g$; multiplying on the right by $g^{-1}$ gives $h_1=h_2$, proving injectivity. Every $x\in Hg$ has the form $hg$ for some $h\in H$, so $x=f(h)$ and the map is surjective. Its inverse is $x\mapsto xg^{-1}$.

\exercise{10.28}{Conjugation condition implies equal cosets}
Suppose $g^{-1}hg\in H$ for all $g\in G,h\in H$. If $hg\in Hg$, write $hg=g(g^{-1}hg)\in gH$, so $Hg\subseteq gH$. For the other inclusion replace $g$ by $g^{-1}$ in the hypothesis: $ghg^{-1}\in H$. Then $gh=(ghg^{-1})g\in Hg$, so $gH\subseteq Hg$. The two inclusions give equality.

\exercise{10.29}{Equal coset partitions imply the conjugation condition}
Assume the partitions agree. The left coset $gH$ and the right coset $Hg$ both contain $g$, so they must be the same cell of that common partition. For $h\in H$, the element $hg$ is in $Hg=gH$, so $hg=gh'$ for some $h'\in H$. Multiplying on the left by $g^{-1}$ gives $g^{-1}hg=h'\in H$, as required.

\exercise{10.30}{Do equal left cosets imply equal right cosets?}
No. In $S_3$, take $H=\{e,(1,2)\}$, $a=(1,3)$, and $b=a(1,2)=(1,2,3)$. Since $b=ah$ for $h\in H$, $aH=bH$. If $Ha=Hb$, the right-coset criterion would give $ba^{-1}\in H$. But
\[ba^{-1}=a(1,2)a^{-1}=(3,2)\notin H.\]
Thus $Ha\ne Hb$, giving a counterexample satisfying the assumed left-coset equality.

\exercise{10.31}{Membership from equal right cosets}
True. Since $e\in H$, the element $b=eb$ belongs to $Hb$. The assumption $Ha=Hb$ then gives $b\in Ha$. This conclusion is a membership statement and does not require commutativity or normality.

\exercise{10.32}{Inverting equal cosets}
True. If $aH=bH$, take the inverses of every element of both sets. Since inversion is a bijection and $H^{-1}=H$,
\[(aH)^{-1}=H^{-1}a^{-1}=Ha^{-1},\qquad (bH)^{-1}=Hb^{-1}.\]
Equal sets have equal inverse sets, so $Ha^{-1}=Hb^{-1}$.

\exercise{10.33}{Squaring representatives of equal cosets}
False. Use $H=\{e,(1,2)\}$, $a=(1,3)$, and $b=(1,2,3)=a(1,2)$ as in Exercise 30. Then $aH=bH$, but $a^2=e$, so $a^2H=H$. Meanwhile $b^2=(1,3,2)\notin H$, so $b^2H\ne H$. Thus squaring a representative need not preserve left-coset equality when the subgroup is not normal.

\exercise{10.34}{Proper subgroups of a group of order $pq$}
By Lagrange, the order of a subgroup divides $pq$. If the subgroup is proper, its order is $1,p$, or $q$ (when $p=q$, the list reduces to $1,p$). The trivial group is cyclic. Any group of prime order is cyclic: a nonidentity element generates a subgroup whose order is a divisor greater than $1$ of that prime, and so generates the whole group. Hence every proper subgroup is cyclic.

\exercise{10.35}{A bijection between left and right cosets}
Define $F$ on left cosets by $F(aH)=Ha^{-1}$. It is well defined because $aH=bH$ implies $Ha^{-1}=Hb^{-1}$, as proved in Exercise 32. Conversely $Ha^{-1}=Hb^{-1}$ implies $aH=bH$ after taking inverse sets, so $F$ is injective. Every right coset $Hg$ equals $F(g^{-1}H)$, proving surjectivity. Thus the collections have the same cardinality even for infinite groups.

\exercise{10.36}{A unique involution in an abelian group of order $2n$}
Existence follows by pairing nonidentity elements with their inverses in a group of even order: at least one is its own inverse. For uniqueness, suppose distinct involutions $a,b$ exist. Since the group is abelian, the four elements $e,a,b,ab$ form a subgroup: all square to $e$, and they are distinct, since $ab=e,a,b$ would force $a=b$, $b=e$, or $a=e$. Thus $4$ divides $2n$ by Lagrange, contradicting oddness of $n$. Hence there is exactly one element of order $2$.

\exercise{10.37}{No proper nontrivial subgroups}
Choose $a\ne e$. Its subgroup $\gen a$ is nontrivial, so by hypothesis it equals $G$ and $G$ is cyclic. It cannot be infinite cyclic, because then $\gen{a^2}$ is a proper nontrivial subgroup: it contains $a^2\ne e$ but not $a$. Thus $G$ has a finite order $m>1$. If $m$ were composite, choose a proper divisor $d$ with $1<d<m$; the element $a^{m/d}$ generates a subgroup of order $d$, again a contradiction. Therefore $m$ is prime.

\exercise{10.38}{Multiplication of subgroup indices}
Let the distinct left cosets of $H$ in $G$ be $a_iH$ for $1\le i\le r$, and those of $K$ in $H$ be $b_jK$ for $1\le j\le s$.
\step{1. Cover $G$.} Given $g\in G$, write $g=a_ih$ for some $i$ and $h\in H$. Then $h=b_jk$ for some $j,k\in K$, so $g\in a_ib_jK$.
\step{2. Show the $rs$ cosets are distinct.} If $a_ib_jK=a_{i'}b_{j'}K$, representatives differ by a factor in $K\subseteq H$, so $a_iH=a_{i'}H$ and $i=i'$. Cancelling $a_i$ then gives $b_jK=b_{j'}K$, hence $j=j'$.
Therefore these are exactly $rs$ distinct left cosets of $K$ in $G$, proving $[G:K]=[G:H][H:K]$. The argument does not require $G$ itself to be finite.

\exercise{10.39}{A subgroup of index two is normal}
If $g\in H$, both $gH$ and $Hg$ equal $H$. If $g\notin H$, the two left cosets are $H$ and $gH$, so $gH=G\setminus H$. The right cosets likewise give $Hg=G\setminus H$. Hence $gH=Hg$ for every $g$, proving normality. In fact the same complement argument works without the given finiteness assumption.

\exercise{10.40}{Every element satisfies $a^n=e$}
Let the order of $a$ be $d$. By Lagrange applied to $\gen a$, $d$ divides $n=|G|$, so write $n=dk$. Then $a^n=(a^d)^k=e^k=e$. The order $d$ is finite because $G$ is finite.

\exercise{10.41}{A unique representative in $[0,1)$}
Let a coset be $a+\Z$. Set $x=a-\lfloor a\rfloor$. Since $\lfloor a\rfloor\le a<\lfloor a\rfloor+1$, we have $0\le x<1$; subtracting the integer $\lfloor a\rfloor$ also shows $x\in a+\Z$. For uniqueness, if $x,y$ both lie in this coset and in $[0,1)$, their difference is an integer with $-1<x-y<1$. The only such integer is $0$, so $x=y$.

\exercise{10.42}{Sine is constant on a coset}
The subgroup generated by $2\pi$ in additive $\R$ is $2\pi\Z$. If $y$ belongs to the coset $x+2\pi\Z$, then $y=x+2\pi k$ for some integer $k$. By the $2\pi$-periodicity of sine, $\sin y=\sin x$. Thus assigning $x+2\pi\Z\mapsto\sin x$ is independent of the representative chosen and defines a function on the coset set.

\exercise{10.43}{Double cosets}
\textbf{(a)} Reflexivity: $a=eae$ with $e\in H,K$. Symmetry: if $a=hbk$, then $b=h^{-1}ak^{-1}$, with the inverse factors still in their subgroups. Transitivity: if $a=h_1bk_1$ and $b=h_2ck_2$, then $a=(h_1h_2)c(k_2k_1)$, with $h_1h_2\in H$ and $k_2k_1\in K$. Thus the relation is an equivalence relation.
\textbf{(b)} Its class containing $a$ consists exactly of all $hak$, with $h\in H$ and $k\in K$, namely the double coset $HaK$. The order of the two factors on the right in the transitivity calculation matters because the group need not be abelian.

\exercise{10.44}{A point stabilizer and its transport coset}
\textbf{(a)} Let $H=S_{c,c}$. The identity fixes $c$; products of maps fixing $c$ fix $c$; and inverses fix $c$ too. Hence $H\le S_A$.
\textbf{(b)} If $d\ne c$, the identity does not send $c$ to $d$, so it is absent from $S_{c,d}$. Therefore this set is not a subgroup.
\textbf{(c)} Choose a permutation $\tau$ sending $c$ to $d$, for example the transposition $(c,d)$. Then $S_{c,d}=\tau H$. Indeed $(\tau h)(c)=\tau(c)=d$ for $h\in H$. Conversely, if $\sigma(c)=d$, then $(\tau^{-1}\sigma)(c)=c$, so $\tau^{-1}\sigma\in H$ and $\sigma\in\tau H$. Thus it is a left coset of the stabilizer.

\exercise{10.45}{The subgroups of a finite cyclic group}
Let $G=\gen g$ have order $n$. For a divisor $d\mid n$, the element $g^{n/d}$ has order $d$, so $H_d=\gen{g^{n/d}}$ is a subgroup of order $d$. Conversely, every subgroup of a cyclic group is cyclic. If $H=\gen{g^k}$ has order $d$, then $d=n/\gcd(n,k)$, so $\gcd(n,k)=n/d$. The subgroup generated by $g^k$ is exactly the subgroup generated by $g^{\gcd(n,k)}$: the exponent gcd is an integer combination of $n,k$, and $g^n=e$. Thus $H=\gen{g^{n/d}}=H_d$. This proves existence and uniqueness for each divisor, and Lagrange excludes other subgroup orders.

\exercise{10.46}{The divisor sum of the Euler function}
Partition the $n$ elements of $G=\gen g$ according to their exact orders. For each divisor $d\mid n$, there is one subgroup of order $d$, and the elements of exact order $d$ are precisely its generators. A cyclic order-$d$ group has $\varphi(d)$ generators, because its generator powers $h^k$ generate exactly when $\gcd(k,d)=1$. Therefore counting all elements gives
\[\boxed{\sum_{d\mid n}\varphi(d)=n.}\]
The divisor $d=1$ contributes the identity, since $\varphi(1)=1$.

\exercise{10.47}{A bound on roots forces cyclicity}
Put $n=|G|$, and for each divisor $d$ of $n$ let $c_d$ count elements of exact order $d$. If $c_d>0$, choose such an element $a$. Its cyclic subgroup has $d$ elements, all satisfying $x^d=e$. By the hypothesis there are at most $d$ solutions in all of $G$, so these are all the solutions. Thus every element of exact order $d$ belongs to $\gen a$, and $c_d=\varphi(d)$. If no such element exists, $c_d=0$. In either case $c_d\le\varphi(d)$.
Lagrange says every element order divides $n$, so
\[n=\sum_{d\mid n}c_d\le\sum_{d\mid n}\varphi(d)=n.\]
Equality of the two sums of nonnegative terms forces equality in every term: otherwise at least one strictly smaller term and no larger terms would make the sum smaller. In particular $c_n=\varphi(n)>0$. An element of order $n$ generates all $n$ elements of $G$, so $G$ is cyclic.
